\medskip

\textbf{Partie A}

\medskip

On considère la suite $\left(u_n\right)_{n\in \N}$ définie par $u_0 = 400$ et pour tout entier naturel $n$ : 
\[u_{n+1} = 0,9u_n + 60.\]

\begin{enumerate}
\item
	\begin{enumerate}
		\item Calculer $u_1$ et $u_2$.
		\item Conjecturer le sens de variation de la suite $\left(u_n\right)_{n\in \N}$
	\end{enumerate}	
\item Montrer, par récurrence, que pour tout entier naturel $n$, on a l'inégalité 
\[0  \leqslant u_n  \leqslant u_{n+1} \leqslant  600.\]
\item 
	\begin{enumerate}
		\item Montrer que la suite $\left(u_n\right)_{n\in \N}$ est convergente.
		\item Déterminer la limite de la suite $\left(u_n\right)_{n\in \N}$. Justifier.
	\end{enumerate}
\item On donne une fonction écrite en langage Python :

\begin{center}
\begin{tabularx}{0.4\linewidth}{|X|}\hline
\textbf{def} mystere(seuil):\\
\qquad n=0\\
\qquad u=400\\
\qquad \textbf{while} u <= seuil :\\
\qquad \qquad n = n+1\\
\qquad \qquad u = 0.9*u+60\\
\qquad \textbf{return} n\\ \hline
\end{tabularx}
\end{center}

Quelle valeur obtient-on en tapant dans la console de Python: mystere (500) ?
\end{enumerate}

\medskip

\textbf{Partie B}

\medskip

Un arboriculteur possède un verger dans lequel il a la place de cultiver au maximum $500$~arbres.

Chaque année il vend 10\,\% des arbres de son verger et puis il replante $60$ nouveaux arbres.

Le verger compte $400$ arbres en 2023.

L'arboriculteur pense qu'il pourra continuer à vendre et à planter les arbres au même rythme pendant les années à venir.

Va-t-il être confronté à un problème de place dans son verger ? Expliquer votre réponse.
